\chapter{朱哥讲金融 /
朱学敏数字金融知识库}\label{ux6731ux54e5ux8bb2ux91d1ux878d--ux6731ux5b66ux654fux6570ux5b57ux91d1ux878dux77e5ux8bc6ux5e93}

\begin{quote}
面向数字货币、数字人民币、CBETS/数币达、稳定币、支付钱包、反洗钱、跨境支付、区块链与
Web3 的主题化知识库。
\end{quote}

\textbf{资料快照：2026-11-16}

\section{使用范围}\label{ux4f7fux7528ux8303ux56f4}

本知识库整理自``朱哥讲金融''\,``我是朱学敏''\,``朱哥讲支付''\,``朱哥聊产品''等公开公众号内容。它保留作者的分析框架、案例和可执行清单，删除重复铺垫、引流文案和宣传性段落。

作者文章属于行业观察与产品经验，不是监管法规原文，也不构成法律或投资建议。涉及监管口径、产品上线状态、牌照、金额和统计数字时，应回到文末列出的官方来源复核。

\section{快速结论}\label{ux5febux901fux7ed3ux8bba}

\begin{enumerate}
\def\labelenumi{\arabic{enumi}.}
\tightlist
\item
  \textbf{数字人民币不是普通第三方支付余额。}
  它是法定货币的数字形态，作者反复用``央行---商业银行/运营机构''双层体系解释其发行和流通。
\item
  \textbf{e-CNY 的产品价值在支付轨道。}
  支付即结算、双离线、可控匿名、钱包分级、智能合约是作者文章中最常复用的底层能力。
\item
  \textbf{CBETS/数币达是跨境网络。}
  它连接境内外参与者，把贸易单据、资金流和信息流放进同一可编程结算流程；其重点不是``快一点''，而是跨境流程重构。
\item
  \textbf{香港稳定币进入持牌与受控互联阶段。} RCHKD
  案例显示，机构可以采用区块链账本形式，同时保留许可制节点、白名单、集中私钥和封闭场景。
\item
  \textbf{风险与效率同源。}
  支付即结算、低门槛钱包和跨域实时转账既提供便利，也带来跑分、拆分交易、快进快出和跨境规避风险。
\item
  \textbf{反洗钱需要贯穿全生命周期。} 核心是分级
  KYC/CDD、钱包限额、名单与制裁筛查、实时行为监测、可疑交易上报、证据保全与用户教育。
\item
  \textbf{Web3 不能与数字人民币混同。}
  前者涉及开放网络、资产代币化和可组合协议；e-CNY
  是国家信用下的法定货币体系。二者可能在技术组件上重叠，但法律属性和治理结构不同。
\end{enumerate}

\section{知识库结构}\label{ux77e5ux8bc6ux5e93ux7ed3ux6784}

\begin{itemize}
\tightlist
\item
  \href{01-核心认知框架.md}{01-核心认知框架}
\item
  \href{02-CBDC与数字人民币.md}{02-CBDC与数字人民币}
\item
  \href{03-CBETS与跨境支付.md}{03-CBETS与跨境支付}
\item
  \href{04-稳定币与香港持牌模式.md}{04-稳定币与香港持牌模式}
\item
  \href{05-反洗钱与风险控制.md}{05-反洗钱与风险控制}
\item
  \href{06-钱包、对公结算与支付运营.md}{06-钱包、对公结算与支付运营}
\item
  \href{07-区块链、Web3与加密金融.md}{07-区块链、Web3与加密金融}
\item
  \href{08-转数快与香港清算.md}{08-转数快与香港清算}
\item
  \href{09-合规检查清单.md}{09-合规检查清单}
\item
  \href{sources/14篇全文来源.md}{sources/14篇全文来源}
\item
  \href{sources/逐篇蒸馏.md}{sources/逐篇蒸馏}
\item
  \href{sources/扩展文章索引.md}{sources/扩展文章索引}
\end{itemize}

\section{来源与方法}\label{ux6765ux6e90ux4e0eux65b9ux6cd5}

\begin{itemize}
\tightlist
\item
  已完整提取并蒸馏 \textbf{14 篇``数字货币''系列文章}，时间覆盖 2026-10-17 至
  2026-11-16。
\item
  通过公开微信搜索索引补充了作者多个账号的公开标题、日期和摘要，用于覆盖稳定币、跨境支付、反洗钱、钱包和
  Web3 等主题。
\item
  公众号页面显示共有 \textbf{122
  篇原创内容}，但公开文章页只能直接展开当前``数字货币''合集；完整账号历史需要微信客户端权限。因此本库不声称已经取得全部
  122 篇全文。
\item
  没有拿到全文的条目只在\href{sources/扩展文章索引.md}{扩展文章索引}中作为线索保留，不把搜索摘要扩写成原文观点。
\item
  对作者文章中容易产生技术误解的表述，已在相应主题页增加``校正''说明。
\end{itemize}

\chapter{核心认知框架}\label{ux6838ux5fc3ux8ba4ux77e5ux6846ux67b6}

\section{一张分层图}\label{ux4e00ux5f20ux5206ux5c42ux56fe}

{\def\LTcaptype{none} % do not increment counter
\begin{longtable}[]{@{}lll@{}}
\toprule\noalign{}
层级 & 解决的问题 & 代表对象 \\
\midrule\noalign{}
\endhead
\bottomrule\noalign{}
\endlastfoot
货币层 & 价值由谁发行、谁的负债、是否法偿 & e-CNY、银行存款、持牌稳定币 \\
网络层 & 付款指令如何到达、何时最终结算 & 央行支付系统、CBETS、转数快 FPS \\
钱包层 & 谁持有、谁控制、额度与权限如何划分 &
个人/对公钱包、母子钱包、软/硬钱包 \\
产品层 & 面向什么场景、怎样交互、怎样退款 &
付款码、NFC、代发工资、红包、跨境贸易 \\
风险层 & 如何识别、阻断、追溯和报告风险 & KYC、CDD、制裁筛查、交易监测、STR \\
数据层 & 交易、凭证、身份和状态怎样记录 &
登记中心、对账台账、许可链、审计日志 \\
\end{longtable}
}

\section{四个不要混淆}\label{ux56dbux4e2aux4e0dux8981ux6df7ux6dc6}

\begin{enumerate}
\def\labelenumi{\arabic{enumi}.}
\tightlist
\item
  \textbf{法定货币与支付机构余额不同。} e-CNY
  是国家信用下的法定货币；微信支付、支付宝等余额是商业机构负债。
\item
  \textbf{区块链与代币不同。}
  许可链可以只承担共享记账和可追溯功能；是否发行可转让代币、是否允许匿名地址参与是另一层设计。
\item
  \textbf{稳定币与 CBDC 不同。} 稳定币通常由持牌发行人负债和储备支持；CBDC
  是中央银行直接或双层体系发行的法定货币。
\item
  \textbf{跨境``更快''不等于自动合规。}
  结算速度提升会压缩人工复核窗口，因此需要前置
  KYC、名单筛查、额度、用途和贸易背景核验。
\end{enumerate}

\section{统一的业务闭环}\label{ux7edfux4e00ux7684ux4e1aux52a1ux95edux73af}

\begin{enumerate}
\def\labelenumi{\arabic{enumi}.}
\tightlist
\item
  身份与钱包开立
\item
  资金来源与用途校验
\item
  支付/转让/兑换
\item
  实时风控与必要的人工复核
\item
  结算最终性与凭证生成
\item
  对账、会计入账与归档
\item
  异常处置、可疑交易报告与复盘
\end{enumerate}

\section{统一的合规闭环}\label{ux7edfux4e00ux7684ux5408ux89c4ux95edux73af}

\texttt{制度\ →\ 客户尽职调查\ →\ 风险评级\ →\ 交易限额\ →\ 名单筛查\ →\ 行为监测\ →\ 可疑报告\ →\ 冻结/退出\ →\ 审计复盘}

任何只强调``支付更快、成本更低''，却没有说明身份、权限、资金用途和异常处置的设计，都不能视为完整方案。

\chapter{CBDC 与数字人民币}\label{cbdc-ux4e0eux6570ux5b57ux4ebaux6c11ux5e01}

\section{定位}\label{ux5b9aux4f4d}

数字人民币（e-CNY）是中国人民银行发行的法定数字货币，主要定位于 M0
数字形态。作者文章反复强调：

\begin{itemize}
\tightlist
\item
  国家信用与法偿性
\item
  央行发行、商业机构参与流通的双层运营
\item
  账户松耦合与钱包载体
\item
  面向零售和对公支付的基础设施属性
\end{itemize}

\section{``一币、两库、三中心''}\label{ux4e00ux5e01ux4e24ux5e93ux4e09ux4e2dux5fc3}

{\def\LTcaptype{none} % do not increment counter
\begin{longtable}[]{@{}ll@{}}
\toprule\noalign{}
组成 & 作用 \\
\midrule\noalign{}
\endhead
\bottomrule\noalign{}
\endlastfoot
一币 & 央行数字货币，即数字人民币 \\
发行库 & 央行侧记录发行基金的数据库 \\
银行库 & 商业银行等运营机构侧的数据库 \\
认证中心 & 管理机构与用户身份信息 \\
登记中心 & 记录权属、交易流水与状态 \\
大数据分析中心 & 反洗钱、支付行为分析与监管指标 \\
\end{longtable}
}

\section{双层运营}\label{ux53ccux5c42ux8fd0ux8425}

\begin{enumerate}
\def\labelenumi{\arabic{enumi}.}
\tightlist
\item
  央行负责发行、规则、基础设施与全生命周期管理。
\item
  指定运营机构面向公众提供兑换和流通服务。
\item
  企业在三层结构中开展业务：央行基础设施 → 运营机构 → 企业/个人钱包。
\end{enumerate}

\section{四个核心特征}\label{ux56dbux4e2aux6838ux5fc3ux7279ux5f81}

\begin{enumerate}
\def\labelenumi{\arabic{enumi}.}
\tightlist
\item
  \textbf{支付即结算}：钱包间价值转移与余额更新同步完成，减少传统清算等待。
\item
  \textbf{双离线支付}：付款端和受理端均离线时仍可先完成交易，联网后再同步。
\item
  \textbf{可控匿名}：小额场景降低身份暴露，大额和可疑交易依法可追溯。
\item
  \textbf{钱包分层}：按主体、载体、权限和 KYC 强度形成不同能力与限额。
\end{enumerate}

\section{钱包分类}\label{ux94b1ux5305ux5206ux7c7b}

\begin{itemize}
\tightlist
\item
  按主体：个人钱包、对公钱包
\item
  按载体：软钱包、硬钱包
\item
  按权限：母钱包、子钱包
\item
  按身份识别强度：不同等级与交易限额
\end{itemize}

\subsection{母子钱包设计}\label{ux6bcdux5b50ux94b1ux5305ux8bbeux8ba1}

母钱包承担强实名、主控和资金归集；子钱包用于部门、项目、员工或外部场景隔离。设计重点不是技术拆分，而是权限、预算、限额和回款路径。

\section{支付场景}\label{ux652fux4ed8ux573aux666f}

\subsection{交通出行}\label{ux4ea4ux901aux51faux884c}

\begin{itemize}
\tightlist
\item
  公交地铁乘车码
\item
  出租车与网约车
\item
  停车缴费、无感支付
\item
  个人钱包向交通运营方对公钱包划转
\item
  退款原路返回
\end{itemize}

\subsection{企业与公共场景}\label{ux4f01ux4e1aux4e0eux516cux5171ux573aux666f}

\begin{itemize}
\tightlist
\item
  代发工资：企业对公钱包批量向员工个人钱包发放
\item
  对公结算：供应链、采购、项目款、费用和税费
\item
  定向支付：智能合约绑定验收、到货、到期或票据状态
\item
  现金红包：现金、群红包和专属红包
\item
  NFC 离线支付：弱网、地下空间、公共交通和应急场景
\end{itemize}

\section{结算与会计}\label{ux7ed3ux7b97ux4e0eux4f1aux8ba1}

企业通常可将数字人民币钱包余额纳入``其他货币资金---数字人民币钱包''科目管理：

\begin{itemize}
\tightlist
\item
  银行账户转入钱包
\item
  钱包向供应商、员工或税务机构付款
\item
  钱包与银行账户互转
\item
  子钱包归集/下拨
\item
  退款、冲正和重复付款处理
\end{itemize}

关键不是会记一笔分录，而是形成``业务依据---审批---支付---实时凭证---日对账---差异闭环---归档''的完整链条。

\section{校正与技术边界}\label{ux6821ux6b63ux4e0eux6280ux672fux8fb9ux754c}

\begin{itemize}
\tightlist
\item
  作者部分文章把``数字人民币''整体描述为基于区块链或分布式账本。更准确的表述是：e-CNY
  采用中心化管理与混合技术架构，区块链/分布式账本可作为部分组件或增强能力，不能因此把
  e-CNY 等同于公链加密资产。
\item
  ``零手续费''属于试点或特定阶段政策口径，不应无条件理解为所有机构、跨境场景和长期商业模式的永久承诺。
\item
  ``双离线''不等于无风险；设备丢失、双花、同步延迟和事后风控都应在产品设计中处理。
\end{itemize}

\section{主要原始文章}\label{ux4e3bux8981ux539fux59cbux6587ux7ae0}

\begin{itemize}
\tightlist
\item
  数币``1币2库3中心4特点''
\item
  数币``付款''
\item
  数币``软钱包''
\item
  数币``交通出行''
\item
  数币``代发工资''
\item
  数币``现金红包''
\item
  数币``NFC支付''
\item
  数币``可控匿名''
\item
  数币结算、对账与会计分录
\item
  数币``对公结算''
\end{itemize}

\chapter{CBETS /
数币达与跨境支付}\label{cbets--ux6570ux5e01ux8fbeux4e0eux8de8ux5883ux652fux4ed8}

\section{定义}\label{ux5b9aux4e49}

CBETS（Cross-border e-CNY Transfer
Services，数币达）是数字人民币跨境基础设施，连接境内外金融机构、支付服务商和终端场景。

核心判断：它不是一个简单的跨境收款按钮，而是``参与者网络 + 规则 + 钱包/账本 +
贸易信息 + 可编程结算''的组合。

\section{与传统跨境支付的区别}\label{ux4e0eux4f20ux7edfux8de8ux5883ux652fux4ed8ux7684ux533aux522b}

传统路径通常经过汇款行、代理行、清算行等多层节点，存在：

\begin{itemize}
\tightlist
\item
  多机构重复审核
\item
  资金状态不透明
\item
  在途时间长
\item
  手续费与汇兑成本叠加
\item
  贸易单据、资金流和信息流分离
\end{itemize}

CBETS
的价值在于统一接入、点对点结算、可编程条件执行和三流协同，而不是只比较``几秒到账''。

\section{三类典型价值}\label{ux4e09ux7c7bux5178ux578bux4ef7ux503c}

\subsection{B2B 贸易结算}\label{b2b-ux8d38ux6613ux7ed3ux7b97}

文章中的老挝电力案例展示了：

\begin{itemize}
\tightlist
\item
  融资款进入指定钱包
\item
  智能合约限定用途
\item
  资金定向支付给供应商
\item
  必要时继续向分包商自动分账
\item
  支付与电子提单等贸易单据挂钩
\end{itemize}

目标是把``专款专用、结算即交付、可追溯''写成可执行规则。

\subsection{跨境电商}\label{ux8de8ux5883ux7535ux5546}

境外买家付款、境内商户收款、汇率与汇回是核心节点。CBETS
可作为连接境外钱包与境内商户的底层网络，缩短跨境链条。

\subsection{B2C 入境消费}\label{b2c-ux5165ux5883ux6d88ux8d39}

境外用户可保留本地银行卡或钱包的充值习惯，在中国境内使用 e-CNY
钱包扫码或碰一碰支付。重点是降低开户、绑卡和身份认证摩擦。

\section{产品设计原则}\label{ux4ea7ux54c1ux8bbeux8ba1ux539fux5219}

\begin{enumerate}
\def\labelenumi{\arabic{enumi}.}
\tightlist
\item
  \textbf{连接器优先}：一次接入触达多个参与机构和国家/地区。
\item
  \textbf{规则可编程}：将付款条件、分账比例和单据校验前置。
\item
  \textbf{资金与信息同步}：交易状态、贸易凭证和会计凭证共同生成。
\item
  \textbf{本地习惯优先}：不要强迫境外用户完全复制中国钱包体验。
\item
  \textbf{异常可逆/可追}：定义错付、退款、撤销、冻结、仲裁和最终性规则。
\end{enumerate}

\section{关键风险}\label{ux5173ux952eux98ceux9669}

\begin{itemize}
\tightlist
\item
  跨境外汇与资本项目管理
\item
  反洗钱、制裁与恐怖融资筛查
\item
  数据跨境、隐私和本地化
\item
  参与者准入、退出与流动性管理
\item
  智能合约漏洞、预言机数据错误和法律可执行性
\item
  运营机构破产、系统中断与结算最终性
\item
  境外监管认可与多法域冲突
\end{itemize}

\section{合规建议}\label{ux5408ux89c4ux5efaux8bae}

\begin{itemize}
\tightlist
\item
  开户前完成参与机构、终端用户和实际受益人尽调。
\item
  对贸易背景、发票、合同、物流和付款路径做一致性校验。
\item
  对高风险国家/地区、异常金额和短期频繁交易设置强化审查。
\item
  将制裁名单、PEP、负面媒体和链上地址情报纳入统一筛查。
\item
  将退款、拒付、冻结、可疑报告和争议处理写入 SLA 与合同。
\end{itemize}

\chapter{稳定币与香港持牌模式}\label{ux7a33ux5b9aux5e01ux4e0eux9999ux6e2fux6301ux724cux6a21ux5f0f}

\section{RCHKD 案例摘要}\label{rchkd-ux6848ux4f8bux6458ux8981}

根据作者 2026-10 的连续两篇文章：

\begin{itemize}
\tightlist
\item
  汇丰于 2026-10-23 在 PayMe 推出港元稳定币 HSBC RedCoin（RCHKD）。
\item
  汇丰持有香港金管局稳定币发行人牌照 FRS02。
\item
  RCHKD 按 1 HKD : 1 RCHKD 设计，由高质量流动性储备资产支持。
\item
  首批场景是 PayMe 内的 P2P 转账与 P2M 商户支付。
\end{itemize}

汇丰官网也确认 RCHKD 为持牌发行、全额储备支持的港元稳定币，可作为事实核验来源。

\section{它不是开放式公链稳定币}\label{ux5b83ux4e0dux662fux5f00ux653eux5f0fux516cux94feux7a33ux5b9aux5e01}

作者对底层结构的拆分最有参考价值：

{\def\LTcaptype{none} % do not increment counter
\begin{longtable}[]{@{}ll@{}}
\toprule\noalign{}
维度 & RCHKD 的设计 \\
\midrule\noalign{}
\endhead
\bottomrule\noalign{}
\endlastfoot
账本 & 许可制 Hyperledger Besu，而非开放公链 \\
接口 & ERC-20 形式兼容 \\
参与 & 节点、地址、用户和渠道受白名单/许可控制 \\
持有 & 客户看到的是份额，代币集中放于 Omnibus 钱包 \\
私钥 & 由汇丰控制，客户取得合同债权 \\
流通 & 限制在汇丰许可的渠道与场景内 \\
可组合性 & 不支持自由接入 DeFi、DEX 或公链桥 \\
\end{longtable}
}

这套设计的核心是``银行负债的链上化''，而不是把一家银行改造成开放加密协议。

\section{发行、转移、赎回}\label{ux53d1ux884cux8f6cux79fbux8d4eux56de}

\begin{itemize}
\tightlist
\item
  发行：把汇丰体系内的港元余额转换为稳定币负债，属于存量负债转换。
\item
  转移：P2P/P2M 在许可渠道内完成，最终以银行渠道账本记录为准。
\item
  赎回：文章描述为按面值兑换港元，目标在一个工作日内完成，但极端赎回、储备变现和法律结构可能产生摩擦。
\end{itemize}

\section{为什么监管偏好这种结构}\label{ux4e3aux4ec0ux4e48ux76d1ux7ba1ux504fux597dux8fd9ux79cdux7ed3ux6784}

\begin{enumerate}
\def\labelenumi{\arabic{enumi}.}
\tightlist
\item
  足额储备与牌照准入提高偿付和治理确定性。
\item
  白名单与 KYC 使反洗钱控制链保持连续。
\item
  渠道隔离便于控制风险传播和跨境边界。
\item
  机构可以复用已有支付用户、商户网络和风控体系。
\end{enumerate}

\section{代价与风险}\label{ux4ee3ux4ef7ux4e0eux98ceux9669}

\begin{itemize}
\tightlist
\item
  放弃无许可网络带来的可组合性。
\item
  用户持有的是对发行人的合同请求权，而非对底层资产的直接物权。
\item
  单机构私钥和集中运营带来操作、治理和集中度风险。
\item
  储备托管、破产隔离、审计频率和赎回压力必须持续核查。
\item
  ``持牌''\,``全额储备''降低风险，但不等于零信用、零操作或零流动性风险。
\end{itemize}

\section{与其他稳定币的比较}\label{ux4e0eux5176ux4ed6ux7a33ux5b9aux5e01ux7684ux6bd4ux8f83}

\begin{itemize}
\tightlist
\item
  与 USDT/USDC：RCHKD 更接近受控银行支付工具，不以开放公链互操作为核心。
\item
  与银行存款：RCHKD 是独立代币负债，可能提供可编程性，但法律保护和清算机制不同。
\item
  与 CBDC：RCHKD 是商业银行负债，不是央行直接发行的法定数字货币。
\end{itemize}

\section{资料来源}\label{ux8d44ux6599ux6765ux6e90}

\begin{itemize}
\tightlist
\item
  作者：《汇丰银行发布港元稳定币RCHKD》
\item
  作者：《扒开``底裤''？研究汇丰银行港元稳定币RCHKD的底层逻辑》
\item
  汇丰官网：\href{https://www.hsbc.com.hk/zh-cn/stablecoins/}{港元稳定币 / 汇丰
  RedCoin}
\end{itemize}

\chapter{反洗钱与风险控制}\label{ux53cdux6d17ux94b1ux4e0eux98ceux9669ux63a7ux5236}

\section{数字支付为什么更脆弱}\label{ux6570ux5b57ux652fux4ed8ux4e3aux4ec0ux4e48ux66f4ux8106ux5f31}

作者重点指出三类高风险机制：

\begin{enumerate}
\def\labelenumi{\arabic{enumi}.}
\tightlist
\item
  \textbf{实时结算压缩拦截窗口}：传统清算延迟无意中提供了人工复核时间；支付即结算会让可疑资金迅速跨账户流转。
\item
  \textbf{低 KYC
  钱包可批量开立}：仅凭手机号开立的小额钱包容易被用于跑分、诈骗和赌博资金中转。
\item
  \textbf{跨区域点对点流转}：脱离多层银行和外汇审核后，资金链路更短，但身份和用途识别更依赖平台侧控制。
\end{enumerate}

\section{常见可疑模式}\label{ux5e38ux89c1ux53efux7591ux6a21ux5f0f}

\begin{itemize}
\tightlist
\item
  多账户短时集中转入、分散转出
\item
  快进快出、夜间高频、金额刻意接近限额
\item
  公转私、私转公用途不明
\item
  异地登录、频繁更换设备
\item
  同一设备控制多钱包
\item
  资金用途与客户经营规模不符
\item
  拆分交易规避单笔或累计限额
\item
  跨境收款来源、贸易单据和物流信息不一致
\end{itemize}

\section{反洗钱控制链}\label{ux53cdux6d17ux94b1ux63a7ux5236ux94fe}

\subsection{客户准入与 CDD}\label{ux5ba2ux6237ux51c6ux5165ux4e0e-cdd}

\begin{itemize}
\tightlist
\item
  实名核验、证件有效性、人脸/活体校验
\item
  实际受益人和实际控制人识别
\item
  业务性质、资金来源、预期交易规模
\item
  高风险国家/地区和行业识别
\item
  风险评级与差异化额度
\end{itemize}

\subsection{KYC 分级钱包}\label{kyc-ux5206ux7ea7ux94b1ux5305}

\begin{itemize}
\tightlist
\item
  低权限钱包仅满足小额日常需求
\item
  高权限钱包强实名、绑定银行卡、留痕更完整
\item
  对余额、单笔、单日和累计交易设置限额
\item
  异常行为触发补认证、降级或冻结
\end{itemize}

\subsection{制裁与名单筛查}\label{ux5236ux88c1ux4e0eux540dux5355ux7b5bux67e5}

\begin{itemize}
\tightlist
\item
  联合国、适用司法辖区制裁名单
\item
  PEP、负面媒体、执法黑名单
\item
  商户、代理行、交易对手和钱包地址
\item
  命中后的复核、拒绝、冻结和申报流程
\end{itemize}

\subsection{实时监测}\label{ux5b9eux65f6ux76d1ux6d4b}

\begin{itemize}
\tightlist
\item
  规则引擎与机器学习结合
\item
  客户、设备、账户、钱包、IP、地理位置的关系图谱
\item
  夜间、异地、高频、拆分和公私互转模型
\item
  评分、预警、限额、冻结和人工复核
\item
  误报率、漏报率和模型漂移持续评估
\end{itemize}

\subsection{可疑交易报告与证据}\label{ux53efux7591ux4ea4ux6613ux62a5ux544aux4e0eux8bc1ux636e}

\begin{itemize}
\tightlist
\item
  标准化 STR 流程
\item
  保存身份、流水、设备、IP、业务单据和模型理由
\item
  必要时冻结、限制和退出客户
\item
  与监管、公安、银行和支付机构协同
\end{itemize}

\section{对稳定币和 Web3
的额外要求}\label{ux5bf9ux7a33ux5b9aux5e01ux548c-web3-ux7684ux989dux5916ux8981ux6c42}

\begin{itemize}
\tightlist
\item
  钱包地址风险评分与资金追踪
\item
  实际控制人和交易对手穿透
\item
  公链/许可链的准入与地址白名单
\item
  托管、私钥和密钥分离
\item
  制裁、被盗资金和混币器暴露筛查
\item
  必要时执行 Travel Rule 等转账信息随行要求
\item
  代币化资产的发行、交易、赎回和二级市场均纳入监测
\end{itemize}

\section{用户侧义务}\label{ux7528ux6237ux4fa7ux4e49ux52a1}

\begin{itemize}
\tightlist
\item
  不出租、出借、出售钱包
\item
  不参与跑分、刷单和代收代付
\item
  不向陌生人提供验证码、生物识别或私钥
\item
  出现异常立即联系运营机构并配合调查
\end{itemize}

\section{硬性判断}\label{ux786cux6027ux5224ux65ad}

反洗钱不是上线后的补丁。钱包分级、交易限额、名单筛查、监测模型、人工复核、报告与审计权限必须在产品设计阶段确定，否则实时支付会同时放大便利和风险。

\chapter{钱包、对公结算与支付运营}\label{ux94b1ux5305ux5bf9ux516cux7ed3ux7b97ux4e0eux652fux4ed8ux8fd0ux8425}

\section{钱包不是简单余额页}\label{ux94b1ux5305ux4e0dux662fux7b80ux5355ux4f59ux989dux9875}

钱包同时承载五项能力：

\begin{enumerate}
\def\labelenumi{\arabic{enumi}.}
\tightlist
\item
  身份与 KYC 等级
\item
  资金持有与隔离
\item
  权限、额度和审批
\item
  交易、流水和凭证
\item
  风控、冻结、退款和客服
\end{enumerate}

\section{母子钱包}\label{ux6bcdux5b50ux94b1ux5305}

\begin{itemize}
\tightlist
\item
  母钱包：强实名、主控、汇总和预算分配
\item
  子钱包：部门、项目、员工、商户或具体用途
\item
  资金归集：子钱包定期上收母钱包
\item
  资金下拨：母钱包按预算和条件向子钱包拨付
\item
  额度控制：按余额、单笔、日累计和用途
\item
  权限隔离：制单、复核、审批、对账互不兼任
\end{itemize}

\section{企业付款标准流程}\label{ux4f01ux4e1aux4ed8ux6b3eux6807ux51c6ux6d41ux7a0b}

\begin{enumerate}
\def\labelenumi{\arabic{enumi}.}
\tightlist
\item
  业务发起并附合同、发票、申请单等依据
\item
  预算与用途校验
\item
  制单：收款方、金额、用途、订单号、子钱包
\item
  自动校验：余额、限额、名单、状态、可疑规则
\item
  双人复核
\item
  大额或特殊交易多级审批
\item
  密码、U 盾、证书或生物识别验证
\item
  支付并生成交易号、回执和凭证
\item
  日/月对账与错账处理
\item
  归档并接受审计
\end{enumerate}

\section{常用支付方式}\label{ux5e38ux7528ux652fux4ed8ux65b9ux5f0f}

\begin{itemize}
\tightlist
\item
  付款码：用户展示，商户扫描
\item
  扫码支付：用户主动扫描商户码
\item
  NFC/碰一碰：弱网或双离线场景
\item
  对公转账：钱包编号、账户或银企直连
\item
  智能合约：到货、验收、到期、开票等条件触发
\end{itemize}

\section{对账}\label{ux5bf9ux8d26}

至少设置四类核对：

\begin{itemize}
\tightlist
\item
  钱包余额与运营机构流水
\item
  业务订单与支付流水
\item
  财务账与银行/钱包账
\item
  汇总账与子钱包/明细账
\end{itemize}

采用``日清日结、差异登记、责任人、处理结果、复核销号''的闭环。

\section{会计处理框架}\label{ux4f1aux8ba1ux5904ux7406ux6846ux67b6}

通常以``其他货币资金---数字人民币钱包''作为企业钱包余额的核算入口：

\begin{itemize}
\tightlist
\item
  银行转入钱包：借钱包，贷银行存款
\item
  收款：借钱包，贷收入/预收账款
\item
  付款：借成本/费用/应付账款，贷钱包
\item
  子钱包归集或下拨：钱包内部科目互转
\item
  错付退回：按原科目冲回并在台账留痕
\end{itemize}

具体科目应按企业会计政策、交易性质、税务处理和审计意见确认。

\section{产品设计要点}\label{ux4ea7ux54c1ux8bbeux8ba1ux8981ux70b9}

\begin{itemize}
\tightlist
\item
  退款链路必须覆盖原路退回、部分退款和跨渠道异常
\item
  离线交易联网后必须自动同步和补对账
\item
  用户看到的是结果，后台必须维护状态机和幂等性
\item
  大额、陌生对手、跨境和公转私应触发强化尽调
\item
  权限变更和离职员工应自动回收钱包权限
\end{itemize}

\chapter{区块链、Web3
与加密金融}\label{ux533aux5757ux94feweb3-ux4e0eux52a0ux5bc6ux91d1ux878d}

\section{作者语料的实际覆盖}\label{ux4f5cux8005ux8bedux6599ux7684ux5b9eux9645ux8986ux76d6}

在本次可访问的文章与搜索索引中，作者对数字人民币、支付、钱包、反洗钱和跨境支付的覆盖明显多于独立
Web3 深度文章。与 Web3
直接相关的可确认条目很少，例如《Obita完成超千万美元天使轮融资》；另有《金融产品方法论》目录提及区块链金融和业务模式。

因此，本页分为两部分：

\begin{itemize}
\tightlist
\item
  \textbf{作者观点}：只保留可以对应到文章标题或正文的内容。
\item
  \textbf{知识补充}：为建立完整框架而补充的基础概念，不冒充作者原文。
\end{itemize}

\section{作者可复用观点}\label{ux4f5cux8005ux53efux590dux7528ux89c2ux70b9}

\begin{itemize}
\tightlist
\item
  金融产品设计应从业务模式、资金流、信息流、权限和风险边界出发，而不是先选技术。
\item
  区块链只有在多方共享状态、可追溯、自动执行或资产可编程确实创造价值时才适合。
\item
  监管合规是金融产品的基础能力；支付、反洗钱、跨境和数据治理不能在上线后补救。
\end{itemize}

\section{基础概念校正}\label{ux57faux7840ux6982ux5ff5ux6821ux6b63}

\subsection{区块链}\label{ux533aux5757ux94fe}

区块链是一种分布式记账技术，通过多方共同维护、密码学签名和共识机制形成可验证的状态记录。它不等于``匿名转账''或``发币''。

\subsection{Web3}\label{web3}

Web3
通常指以开放网络、用户自持身份/资产、智能合约和可组合协议为核心的互联网形态。其具体含义依技术、投资和监管语境而不同，不能把``Web3''当成合规或技术安全的证明。

\subsection{加密金融}\label{ux52a0ux5bc6ux91d1ux878d}

加密金融是围绕加密资产、链上协议、托管、交易、借贷、支付和衍生服务形成的业务集合。法律属性因司法辖区而异。

\section{三种可以落地的链上金融模式}\label{ux4e09ux79cdux53efux4ee5ux843dux5730ux7684ux94feux4e0aux91d1ux878dux6a21ux5f0f}

\subsection{许可链支付与结算}\label{ux8bb8ux53efux94feux652fux4ed8ux4e0eux7ed3ux7b97}

\begin{itemize}
\tightlist
\item
  参与节点和地址白名单
\item
  KYC 与反洗钱控制连续
\item
  共享账本、可追溯和自动对账
\item
  适用于银行、支付机构和供应链联盟
\end{itemize}

\subsection{代币化资产与 RWA}\label{ux4ee3ux5e01ux5316ux8d44ux4ea7ux4e0e-rwa}

\begin{itemize}
\tightlist
\item
  债券、基金份额、应收账款、商品和房地产权益代币化
\item
  可分裂、可组合、支持自动付息与条件转让
\item
  核心难点在法律权利映射、托管、估值、流动性和二级市场合规
\end{itemize}

\subsection{稳定币结算}\label{ux7a33ux5b9aux5e01ux7ed3ux7b97}

\begin{itemize}
\tightlist
\item
  以法币或高质量流动资产支持
\item
  可支持 P2P、P2M、跨境支付和投资品认购
\item
  发行主体信用、储备透明度和监管牌照决定风险等级
\end{itemize}

\section{加密金融合规矩阵}\label{ux52a0ux5bc6ux91d1ux878dux5408ux89c4ux77e9ux9635}

{\def\LTcaptype{none} % do not increment counter
\begin{longtable}[]{@{}lll@{}}
\toprule\noalign{}
风险面 & 核心问题 & 必要控制 \\
\midrule\noalign{}
\endhead
\bottomrule\noalign{}
\endlastfoot
发行 & 是否构成证券、支付工具或稳定币 & 法律定性、牌照、白皮书/披露 \\
托管 & 客户资产是否真正隔离 & 托管牌照、私钥管理、审计 \\
交易 & 是否提供撮合、兑换或杠杆 & 交易所/经纪/衍生品许可 \\
洗钱 & 地址是否受制裁或涉非法资金 & 钱包评分、KYT、Travel Rule \\
市场行为 & 是否存在操纵、内幕和误导宣传 & 监测、广告审核、投诉机制 \\
数据 & 链上信息是否涉及个人数据 & 最小化、跨境与隐私评估 \\
技术 & 合约漏洞、预言机和桥风险 & 审计、监控、限额、应急停机 \\
跨境 & 不同法域规则冲突 & 地域限制、法律意见、牌照映射 \\
\end{longtable}
}

\section{中国大陆与香港的边界}\label{ux4e2dux56fdux5927ux9646ux4e0eux9999ux6e2fux7684ux8fb9ux754c}

\begin{itemize}
\tightlist
\item
  中国大陆对虚拟货币相关非法金融活动的监管框架与数字人民币试点不能混为一谈。e-CNY
  是法定货币，不是加密资产。
\item
  香港允许在牌照与监管框架下开展部分虚拟资产、稳定币和代币化业务，但具体活动仍需按业务实质判断。
\item
  做跨境产品时，应分别判断发行地、用户所在地、资产所在地、服务器所在地和支付清算路径。
\end{itemize}

\section{对外部资料的核验入口}\label{ux5bf9ux5916ux90e8ux8d44ux6599ux7684ux6838ux9a8cux5165ux53e3}

\begin{itemize}
\tightlist
\item
  中国人民银行数字人民币研发工作组：《中国数字人民币的研发进展白皮书》
\item
  香港金管局：稳定币、虚拟资产及转数快相关监管与产品资料
\item
  汇丰银行 RCHKD 官方页面
\item
  项目方披露与独立审计报告
\end{itemize}

\chapter{转数快（FPS）与香港清算}\label{ux8f6cux6570ux5febfpsux4e0eux9999ux6e2fux6e05ux7b97}

\section{基本定位}\label{ux57faux672cux5b9aux4f4d}

转数快（Faster Payment System，FPS）是香港金管局于 2018
年推出的快速支付系统，由香港银行同业结算有限公司运作，连接银行与储值支付工具营运商，支持全天候本地即时转账和支付。

\section{产品能力}\label{ux4ea7ux54c1ux80fdux529b}

\begin{itemize}
\tightlist
\item
  支持港元及人民币（按参与机构能力）
\item
  可通过手机号、电邮、FPS 识别码或二维码收款
\item
  支持 P2P 转账、商户支付和钱包增值
\item
  跨银行、跨电子钱包
\item
  7×24 小时运行
\end{itemize}

\section{在 e-CNY
与跨境支付中的位置}\label{ux5728-e-cny-ux4e0eux8de8ux5883ux652fux4ed8ux4e2dux7684ux4f4dux7f6e}

公开监管与媒体资料显示，e-CNY 曾探索与 FPS
互通，以支持香港用户更便利地充值和使用数字人民币。两者可以形成互补：

\begin{itemize}
\tightlist
\item
  FPS：香港本地银行与钱包之间的清算型基础设施
\item
  e-CNY：法定数字人民币的钱包与支付体系
\item
  CBETS：e-CNY 跨境参与者与场景网络
\end{itemize}

三者不能被简单描述为同一个系统，也不能默认任何组合都已在全国或所有机构正式商用。

\section{风控重点}\label{ux98ceux63a7ux91cdux70b9}

\begin{itemize}
\tightlist
\item
  手机号、邮箱和 FPS ID 与实际收款人姓名的绑定校验
\item
  首笔大额、陌生收款人和名称不一致预警
\item
  诈骗账户、钱骡账户和异常设备识别
\item
  单笔与日累计限额
\item
  付款前展示收款机构与核验名称
\item
  争议、误转、撤回和司法调证流程
\end{itemize}

\section{作者语料覆盖说明}\label{ux4f5cux8005ux8bedux6599ux8986ux76d6ux8bf4ux660e}

本次可通过公开索引访问的``数字货币''系列中没有发现一篇以``转数快''为标题的完整作者文章。因此本页是基础补充页，不作为作者原文摘要。

\section{官方与公开资料}\label{ux5b98ux65b9ux4e0eux516cux5f00ux8d44ux6599}

\begin{itemize}
\tightlist
\item
  \href{https://fps.hkicl.com.hk/schi/fps/index.php}{HKICL：转数快}
\item
  \href{https://www.hkma.gov.hk/}{香港金管局：转数快介绍}
\end{itemize}

\chapter{数字支付与加密金融合规检查清单}\label{ux6570ux5b57ux652fux4ed8ux4e0eux52a0ux5bc6ux91d1ux878dux5408ux89c4ux68c0ux67e5ux6e05ux5355}

\section{业务定性}\label{ux4e1aux52a1ux5b9aux6027}

\begin{itemize}
\tightlist
\item
  产品是否构成支付、存款、稳定币、证券、基金、资管或衍生品？
\item
  发行人、运营者、托管人和商户分别承担什么法律角色？
\item
  是否涉及人民币、港元、外币或加密资产？
\item
  用户、资产、服务器和清算路径分别位于哪些司法辖区？
\end{itemize}

\section{牌照与准入}\label{ux724cux7167ux4e0eux51c6ux5165}

\begin{itemize}
\tightlist
\item
  是否需要支付、银行、电子货币、稳定币、虚拟资产或证券牌照？
\item
  合作方是否持牌，牌照是否覆盖拟开展活动？
\item
  外包、代理和导流是否改变业务实质？
\item
  境外牌照能否覆盖当地用户？
\end{itemize}

\section{客户与钱包}\label{ux5ba2ux6237ux4e0eux94b1ux5305}

\begin{itemize}
\tightlist
\item
  开户身份信息、实控人和实际受益人是否完整？
\item
  钱包等级、余额和交易限额是否与风险相匹配？
\item
  母子钱包权限是否隔离？
\item
  设备、登录、密码、生物识别和私钥是否安全？
\item
  钱包冻结、降级、注销和继承流程是否明确？
\end{itemize}

\section{资金与资产隔离}\label{ux8d44ux91d1ux4e0eux8d44ux4ea7ux9694ux79bb}

\begin{itemize}
\tightlist
\item
  客户资金/资产是否与机构自有资产隔离？
\item
  储备资产范围、托管、审计和披露频率是否明确？
\item
  稳定币赎回权是物权还是合同债权？
\item
  破产隔离和极端赎回情形是否有法律意见？
\end{itemize}

\section{反洗钱与制裁}\label{ux53cdux6d17ux94b1ux4e0eux5236ux88c1}

\begin{itemize}
\tightlist
\item
  是否执行 KYC/CDD、EDD、PEP 和实际受益人识别？
\item
  是否进行制裁、负面媒体和钱包地址筛查？
\item
  是否监测拆分、快进快出、异常跨境和公转私？
\item
  是否具备 STR、冻结、退出和证据留存流程？
\item
  是否进行模型验证、误报分析和审计跟踪？
\end{itemize}

\section{交易与市场合规}\label{ux4ea4ux6613ux4e0eux5e02ux573aux5408ux89c4}

\begin{itemize}
\tightlist
\item
  是否有清晰费率、汇率、退款和争议机制？
\item
  广告是否存在收益承诺、误导或伪装成银行/法币的表述？
\item
  是否存在杠杆、借贷、撮合、做市或衍生功能？
\item
  是否具备市场操纵和异常交易监测？
\end{itemize}

\section{数据与网络安全}\label{ux6570ux636eux4e0eux7f51ux7edcux5b89ux5168}

\begin{itemize}
\tightlist
\item
  数据收集是否最小化？
\item
  数据跨境传输是否有合法依据和评估？
\item
  是否实施密钥分离、多人审批、灾备和渗透测试？
\item
  智能合约是否独立审计，是否有暂停和升级权限？
\end{itemize}

\section{财务与审计}\label{ux8d22ux52a1ux4e0eux5ba1ux8ba1}

\begin{itemize}
\tightlist
\item
  钱包余额、储备和客户权益是否每日核对？
\item
  业务订单、支付流水、财务账和链上记录能否四流一致？
\item
  退款、冲正、错付和跨日交易是否有完整分录？
\item
  审计人员能否取得独立证据？
\end{itemize}

\section{消费者保护}\label{ux6d88ux8d39ux8005ux4fddux62a4}

\begin{itemize}
\tightlist
\item
  用户是否理解持币人权利、发行人信用和赎回限制？
\item
  是否明确不可逆交易、错误地址和私钥丢失后果？
\item
  是否设有投诉、冻结、申诉和赔付机制？
\item
  是否进行适老化、反诈和风险教育？
\end{itemize}

\section{上线决策}\label{ux4e0aux7ebfux51b3ux7b56}

以下任一项为``否''，不应直接上线：

\begin{itemize}
\tightlist
\item
  法律定性和牌照要求已明确
\item
  客户资金/储备隔离可验证
\item
  来源与用途可追溯
\item
  黑名单和可疑交易可阻断
\item
  结算最终性和异常处理可执行
\item
  对账、会计与审计闭环成立
\item
  隐私和跨境数据路径合规
\end{itemize}

\chapter{逐篇蒸馏}\label{ux9010ux7bc7ux84b8ux998f}

只总结全文已实际获取的 14
篇文章。标题、日期和链接见\href{14篇全文来源.md}{14篇全文来源}。

\section{数币``交通出行''}\label{ux6570ux5e01ux4ea4ux901aux51faux884c}

\begin{itemize}
\tightlist
\item
  数字人民币可用于公交地铁乘车码、出租车/网约车和停车缴费。
\item
  资金路径是个人钱包向交通或停车运营方的对公钱包划转，退款原路返回。
\item
  弱网、双离线、可追溯和低成本是其在交通场景的主要优势。
\item
  产品重点在支付状态、退款、清分和运营方对账，不只是生成一个二维码。
\end{itemize}

\section{数币``反洗钱''}\label{ux6570ux5e01ux53cdux6d17ux94b1}

\begin{itemize}
\tightlist
\item
  支付即结算、弱实名钱包、跨区域点对点流转会提高资金追溯和拦截难度。
\item
  典型风险包括跑分、赌博、诈骗、批次小额钱包、多账户拆分和快速分散。
\item
  控制体系应覆盖 KYC/CDD、分级钱包、限额、智能监测、STR
  可疑交易报告和全链路证据。
\item
  运营机构、监管部门和用户都需要承担相应责任，不能只依赖事后冻结。
\end{itemize}

\section{数币``可控匿名''}\label{ux6570ux5e01ux53efux63a7ux533fux540d}

\begin{itemize}
\tightlist
\item
  可控匿名是在合理隐私与可依法追溯之间取平衡。
\item
  小额匿名钱包满足日常支付；高权限、大额交易需要更强实名和审查。
\item
  子钱包推送和 ID 匿名化可减少商业平台接触用户真实身份。
\item
  匿名不是无监管，大额或可疑交易仍需依法穿透。
\end{itemize}

\section{数币``代发工资''}\label{ux6570ux5e01ux4ee3ux53d1ux5de5ux8d44}

\begin{itemize}
\tightlist
\item
  企业和单位需要开通对公钱包，员工需要个人实名钱包。
\item
  企业收集钱包编号，准备资金后按批量文件发放并获取结果文件。
\item
  文章强调实时到账、可追溯、降低手续费用和条件化智能合约。
\item
  实操必须处理员工身份、钱包编号校验、重复付款、退回和工资数据保护。
\end{itemize}

\section{数币``付款''}\label{ux6570ux5e01ux4ed8ux6b3e}

\begin{itemize}
\tightlist
\item
  主要方式包括付款码、主动扫码、NFC 碰一碰和对公转账/定向支付。
\item
  标准流程是身份与对方校验、余额检查、密码/生物识别验证、支付、回执和记账。
\item
  企业付款应实行制单、复核、审批、对账岗位分离，并保留合同、发票等业务依据。
\item
  高频、夜间、陌生对手、异常大额、拆分和公私互转需要重点监控。
\end{itemize}

\section{数币``NFC支付''}\label{ux6570ux5e01nfcux652fux4ed8}

\begin{itemize}
\tightlist
\item
  NFC 利用近距离通信实现设备间点对点支付。
\item
  在弱网、地下室、公共交通等环境，离线 NFC 提升了可用性。
\item
  文章强调交易签名、设备绑定、生物识别和手机丢失时的防护。
\item
  产品仍需定义双离线数据同步、限额、双花防护和离线交易撤销。
\end{itemize}

\section{汇丰银行发布港元稳定币RCHKD}\label{ux6c47ux4e30ux94f6ux884cux53d1ux5e03ux6e2fux5143ux7a33ux5b9aux5e01rchkd}

\begin{itemize}
\tightlist
\item
  文章称汇丰于 2026-10-23 在 PayMe 上线 RCHKD，按港元 1:1 设计。
\item
  首批场景为 P2P 转账和 P2M 商户支付。
\item
  作者将其定位为传统移动支付与持牌稳定币结合的零售支付试验。
\item
  官方产品页可确认 RCHKD 与 FRS02 牌照，但业务数据仍应以发行方最新披露为准。
\end{itemize}

\section{扒开``底裤''？研究汇丰银行港元稳定币RCHKD的底层逻辑}\label{ux6252ux5f00ux5e95ux88e4ux7814ux7a76ux6c47ux4e30ux94f6ux884cux6e2fux5143ux7a33ux5b9aux5e01rchkdux7684ux5e95ux5c42ux903bux8f91}

\begin{itemize}
\tightlist
\item
  RCHKD 使用许可制 Hyperledger Besu 与 ERC-20 接口，不是开放公链资产。
\item
  参与者、节点和地址受白名单控制；客户代币集中于 Omnibus 钱包，私钥由汇丰控制。
\item
  客户持有更接近对发行人的合同请求权，渠道之间隔离，不能自由接入 DeFi。
\item
  其优势是合规确定性、银行信用和支付网络；代价是开放可组合性、跨渠道自由和去中心化。
\end{itemize}

\section{数币``现金红包''}\label{ux6570ux5e01ux73b0ux91d1ux7ea2ux5305}

\begin{itemize}
\tightlist
\item
  红包分为现金红包、群红包和专属红包。
\item
  群红包支持拼手气、拼手速和平均金额等形式。
\item
  资金最终仍是 e-CNY，社交平台承担传播和触达渠道，不改变货币属性。
\item
  文中增长速度和跨平台覆盖等数字缺少独立核验，不应作为当前市场事实直接引用。
\end{itemize}

\section{数币结算、对账与会计分录}\label{ux6570ux5e01ux7ed3ux7b97ux5bf9ux8d26ux4e0eux4f1aux8ba1ux5206ux5f55}

\begin{itemize}
\tightlist
\item
  企业结算流程应从业务依据、制单、系统校验、双人复核、审批、支付、回执到归档。
\item
  大额、陌生对手、公转私、跨境或保证金交易需要更高等级审批。
\item
  对账应覆盖钱包余额、运营机构流水、业务台账和财务账，并设置差异调节表。
\item
  会计上通常围绕``其他货币资金---数字人民币钱包''记录银行互转、收付款、退回、归集和智能合约支付。
\end{itemize}

\section{数币``1币2库3中心4特点''}\label{ux6570ux5e011ux5e012ux5e933ux4e2dux5fc34ux7279ux70b9}

\begin{itemize}
\tightlist
\item
  一币指数字人民币；两库指发行库与银行库。
\item
  三中心分别负责认证、权属/交易登记和大数据分析。
\item
  双层运营由央行发行、商业机构面向用户兑换和流通。
\item
  四项特征被作者总结为支付即结算、双离线、可控匿名和钱包多样化。
\end{itemize}

\section{CBETS数币达}\label{cbetsux6570ux5e01ux8fbe}

\begin{itemize}
\tightlist
\item
  CBETS 是连接境内外金融机构、支付服务商和用户的跨境基础设施。
\item
  B2B 贸易可利用智能合约实现融资定向支付、自动分账和三流合一。
\item
  B2C 场景重点在境外用户入境后保持原有充值习惯，同时使用境内 e-CNY 支付。
\item
  其产品启示是``一次接入、连接全网''，但跨境合规、流动性和结算最终性是关键约束。
\end{itemize}

\section{数币``软钱包''}\label{ux6570ux5e01ux8f6fux94b1ux5305}

\begin{itemize}
\tightlist
\item
  软钱包依靠非对称加密、密码、手势和生物识别保护访问。
\item
  钱包与银行账户可绑定，但通过子钱包和分级权限维持独立的支付入口。
\item
  文章把分布式账本、可追溯和离线支付列为软钱包的技术支撑。
\item
  技术校正：不能据此把 e-CNY
  等同于公有链加密资产；需区分央行中心化治理与部分分布式技术组件。
\end{itemize}

\section{数币``对公结算''}\label{ux6570ux5e01ux5bf9ux516cux7ed3ux7b97}

\begin{itemize}
\tightlist
\item
  e-CNY 对公结算的核心是实时、可追溯、可编程和全天候资金划转。
\item
  企业应建立主/子钱包、审批、限额、业务依据、凭证和对账完整闭环。
\item
  智能合约用于专款专用、到货/验收付款、自动分账和补贴发放。
\item
  方案落地必须处理账户权限、资金归集、退款、错付、会计政策和反洗钱义务。
\end{itemize}

\chapter{14 篇全文来源}\label{14-ux7bc7ux5168ux6587ux6765ux6e90}

以下文章已从公众号正文提取，并在主题页中做了压缩整理。按发布时间排序。

{\def\LTcaptype{none} % do not increment counter
\begin{longtable}[]{@{}>{\raggedright\arraybackslash}p{1.9cm}>{\raggedright\arraybackslash}p{2.8cm}>{\raggedright\arraybackslash}p{6.0cm}>{\raggedright\arraybackslash}p{4.7cm}@{}}
\toprule\noalign{}
日期 & 标题 & 公众号 & 原文 \\
\midrule\noalign{}
\endhead
\bottomrule\noalign{}
\endlastfoot
2026-10-17 21:32 & 数币``交通出行'' & 朱哥讲金融 &
\href{https://mp.weixin.qq.com/s?__biz=MzkyMjE4MDgxOA==&mid=2247508403&idx=1&sn=b3b7ca12804bdbec2d81ddf72c0ee365&chksm=c1faac0ef68d25188533598bc5e8f99c1ec871defe042fd22ced8db8629578da378c8986e3ef&cur_album_id=4742095690531438593&scene=189}{原文} \\
2026-10-18 19:48 & 数币``反洗钱'' & 朱哥讲金融 &
\href{https://mp.weixin.qq.com/s?__biz=MzkyMjE4MDgxOA==&mid=2247508425&idx=1&sn=4cb3cf2df92d703cc01f52e90d1b699a&chksm=c1faac74f68d2562e5634e16e9c7405820bf1d0bdfaedea79702e87f63b270779091dfe7792a&cur_album_id=4742095690531438593&scene=189}{原文} \\
2026-10-21 09:11 & 数币``可控匿名'' & 朱哥讲金融 &
\href{https://mp.weixin.qq.com/s?__biz=MzkyMjE4MDgxOA==&mid=2247508472&idx=1&sn=0105bebb84da123352ffaac4a7a9bdf5&chksm=c1faac45f68d25535d1a46e6280d3578dab28c23bf9711130648e4694207ba6e91a7051f8fad&cur_album_id=4742095690531438593&scene=189}{原文} \\
2026-10-23 09:21 & 数币``代发工资'' & 朱哥讲金融 &
\href{https://mp.weixin.qq.com/s?__biz=MzkyMjE4MDgxOA==&mid=2247508482&idx=1&sn=0960caf789029297fce1e69c8f76a282&chksm=c1faafbff68d26a9b84fe90c15b1c1d5110e2490e3b0cba09365070d1af869e19bd85c4c5c77&cur_album_id=4742095690531438593&scene=189}{原文} \\
2026-10-24 01:10 & 数币``付款'' & 朱哥讲金融 &
\href{https://mp.weixin.qq.com/s?__biz=MzkyMjE4MDgxOA==&mid=2247508492&idx=1&sn=d43ce7ad328f6a757ad5db91bb028c49&chksm=c1faafb1f68d26a727cb38b98aaf3f846742e2d87621a7b2dd51086a7b9fbcfa131a26cf4910&cur_album_id=4742095690531438593&scene=189}{原文} \\
2026-10-24 10:59 & 数币``NFC支付'' & 朱哥讲金融 &
\href{https://mp.weixin.qq.com/s?__biz=MzkyMjE4MDgxOA==&mid=2247508497&idx=1&sn=be76428424f93bc55f8ad6fa441061fd&chksm=c1faafacf68d26bab832629e24644db441ee026433a586b1390d6b1229115444b54dbab6c2e4&cur_album_id=4742095690531438593&scene=189}{原文} \\
2026-10-24 22:50 & 汇丰银行发布港元稳定币RCHKD & 朱哥讲金融 &
\href{https://mp.weixin.qq.com/s?__biz=MzkyMjE4MDgxOA==&mid=2247508498&idx=1&sn=14fae70c44229ca9625dcda14b0f50a1&chksm=c1faafaff68d26b9ba34587169559f812c5acf84a18f1ab34aa4513ab11dad214bc81ffa17b2&cur_album_id=4742095690531438593&scene=189}{原文} \\
2026-10-26 13:28 & 扒开``底裤''？研究汇丰银行港元稳定币RCHKD的底层逻辑 &
朱哥讲金融 &
\href{https://mp.weixin.qq.com/s?__biz=MzkyMjE4MDgxOA==&mid=2247508500&idx=1&sn=55ffe814ab132366ca92dd263ee1f90b&chksm=c1faafa9f68d26bf3201bf93d96bce8c60696cc8c719f4b12c6b56e7fe5daa990a8c134a57e5&cur_album_id=4742095690531438593&scene=189}{原文} \\
2026-10-27 16:44 & 数币``现金红包'' & 朱哥讲金融 &
\href{https://mp.weixin.qq.com/s?__biz=MzkyMjE4MDgxOA==&mid=2247508507&idx=1&sn=b5b49204af3dc30ab048d4332be46d17&chksm=c1faafa6f68d26b02258e35a212130d8ab13744370c2f3dbc65fe3009fa33703df227954c546&cur_album_id=4742095690531438593&scene=189}{原文} \\
2026-10-31 15:46 & 数币结算、对账与会计分录 & 朱哥讲金融 &
\href{https://mp.weixin.qq.com/s?__biz=MzkyMjE4MDgxOA==&mid=2247508512&idx=1&sn=6fda80c7a9ab88dd3c405dc74229209a&chksm=c1faaf9df68d268b281d3d6701ec9255f87fe3d6999b185ef7468ef26f20b482a85a50b0ad47&cur_album_id=4742095690531438593&scene=189}{原文} \\
2026-11-04 18:47 & 数币``软钱包'' & 朱哥讲金融 &
\href{https://mp.weixin.qq.com/s?__biz=MzkyMjE4MDgxOA==&mid=2247508522&idx=1&sn=b7cf30b039e5950d6d1a3716b9f56518&chksm=c1faaf97f68d2681002364f259378137caeb0190e7bdf8f194f7b48efcd946515ea046486a56&cur_album_id=4742095690531438593&scene=189}{原文} \\
2026-11-06 00:04 & 数币``对公结算'' & 朱哥讲金融 &
\href{https://mp.weixin.qq.com/s?__biz=MzkyMjE4MDgxOA==&mid=2247508529&idx=2&sn=96462ea21a125d56fffe27b0045ffab9&chksm=c1faaf8cf68d269ad0b9a5ee33bbf32a177067a02d06b30008824f7578c28d0d798d6a8be2c9&cur_album_id=4742095690531438593&scene=189}{原文} \\
2026-11-10 23:08 & 数币``1币2库3中心4特点'' & 朱哥讲金融 &
\href{https://mp.weixin.qq.com/s?__biz=MzkyMjE4MDgxOA==&mid=2247508572&idx=1&sn=f03c6c708abaae09698c01c615a8b74a&chksm=c1faafe1f68d26f7b69e9171a6fe947601e977aafc44e05e40bb23f7990e71e48b507bc7074d&cur_album_id=4742095690531438593&scene=189}{原文} \\
2026-11-16 21:09 & CBETS数币达 & 朱哥讲金融 &
\href{https://mp.weixin.qq.com/s/1g2YASC79NqbWAOnHN5NFg}{原文} \\
\end{longtable}
}

\chapter{扩展文章索引}\label{ux6269ux5c55ux6587ux7ae0ux7d22ux5f15}

本页收录通过公开微信搜索索引发现的作者相关标题。除已列入``14
篇全文来源''的文章外，其他条目只拿到标题、日期和搜索摘要，不能视为已经完整读取全文。

检索口径：\texttt{朱哥讲金融}、\texttt{朱学敏}、\texttt{我是朱学敏}、\texttt{朱哥讲支付}、\texttt{朱哥聊产品}，叠加数字货币、CBDC、稳定币、区块链、数币达、转数快、反洗钱、跨境支付、Web3、钱包、二清合规等关键词。

\section{数字货币与数字人民币}\label{ux6570ux5b57ux8d27ux5e01ux4e0eux6570ux5b57ux4ebaux6c11ux5e01}

{\def\LTcaptype{none} % do not increment counter
\begin{longtable}[]{@{}>{\raggedright\arraybackslash}p{1.9cm}>{\raggedright\arraybackslash}p{2.8cm}>{\raggedright\arraybackslash}p{6.0cm}>{\raggedright\arraybackslash}p{4.7cm}@{}}
\toprule\noalign{}
日期 & 公众号 & 标题 & 可用信息 \\
\midrule\noalign{}
\endhead
\bottomrule\noalign{}
\endlastfoot
2026-04-01 & 朱哥讲支付 & 数字人民币赋能供应链金融创新 &
数字人民币用于供应链融资、资金定向支付和智能合约 \\
2026-03-21 & 朱哥讲金融 & 银行``母子钱包'' &
母钱包强实名，子钱包用于场景隔离、额度与权限控制 \\
2026-11-06 & 我是朱学敏 & 企业钱包管理 &
企业钱包内控覆盖反洗钱、欺诈、越权、挪用、信息泄露与账户冒用 \\
2026-03-26 & 朱哥讲金融 & 交通银行BOCOM &
与数字人民币、金融服务生态相关，需全文复核 \\
\end{longtable}
}

\section{稳定币与数字资产}\label{ux7a33ux5b9aux5e01ux4e0eux6570ux5b57ux8d44ux4ea7}

{\def\LTcaptype{none} % do not increment counter
\begin{longtable}[]{@{}>{\raggedright\arraybackslash}p{1.9cm}>{\raggedright\arraybackslash}p{2.8cm}>{\raggedright\arraybackslash}p{6.0cm}>{\raggedright\arraybackslash}p{4.7cm}@{}}
\toprule\noalign{}
日期 & 公众号 & 标题 & 可用信息 \\
\midrule\noalign{}
\endhead
\bottomrule\noalign{}
\endlastfoot
2026-10-24 & 朱哥讲金融 & 汇丰银行发布港元稳定币RCHKD & 已取得全文并整理 \\
2026-10-26 & 朱哥讲金融 & 扒开``底裤''？研究汇丰银行港元稳定币RCHKD的底层逻辑 &
已取得全文并整理 \\
近期 & 我是朱学敏 & Obita完成超千万美元天使轮融资 &
搜索摘要只显示融资标题，需全文复核其 Web3 业务细节 \\
\end{longtable}
}

\section{反洗钱、风控与支付合规}\label{ux53cdux6d17ux94b1ux98ceux63a7ux4e0eux652fux4ed8ux5408ux89c4}

{\def\LTcaptype{none} % do not increment counter
\begin{longtable}[]{@{}>{\raggedright\arraybackslash}p{1.9cm}>{\raggedright\arraybackslash}p{2.8cm}>{\raggedright\arraybackslash}p{6.0cm}>{\raggedright\arraybackslash}p{4.7cm}@{}}
\toprule\noalign{}
日期 & 公众号 & 标题 & 可用信息 \\
\midrule\noalign{}
\endhead
\bottomrule\noalign{}
\endlastfoot
2026-03-22 & 朱哥讲金融 & 强监管下合规成生存底线，2026年支付行业罚单金额逼近亿元
& 反洗钱履职、客户身份识别和大额/可疑交易报送是处罚重点 \\
2026-03-27 & 朱哥讲金融 & 跨境支付风控量化 &
反洗钱与制裁筛查应覆盖支付前、中、后各阶段 \\
2026-04-01 & 朱哥讲支付 & 跨境支付``前中后''风控量化 &
同上，标题索引存在版本差异 \\
2026-02-26 & 朱哥讲金融 & 2026年支付行业首张千万级罚单，直指银盛支付违规乱象 &
反洗钱失职、外包管理失责与``双罚制'' \\
2026-02-28 & 朱哥讲金融 & 银行业罚单超1060张，累计罚没超8.55亿 &
监管处罚与银行内控主题 \\
近一年 & 朱哥讲支付 & 非银行支付机构监督管理条例 &
非银支付牌照、业务边界与合规要求 \\
近一年 & 朱哥讲支付 & 银行卡要素验证2/3/4/5/9 &
身份核验、二三四要素等支付风控方法 \\
近一年 & 华创微课金融企服 & 搞懂支付二清合规方案 &
聚合支付、资金池与二清风险隔离 \\
近一年 & 朱哥讲支付 & 支付二清合规解决方案 & 同上 \\
近一年 & 朱哥讲支付 & 一文搞懂支付反洗钱AML & AML/KYC、监测与报告主题 \\
\end{longtable}
}

\section{跨境支付与清算}\label{ux8de8ux5883ux652fux4ed8ux4e0eux6e05ux7b97}

{\def\LTcaptype{none} % do not increment counter
\begin{longtable}[]{@{}>{\raggedright\arraybackslash}p{1.9cm}>{\raggedright\arraybackslash}p{2.8cm}>{\raggedright\arraybackslash}p{6.0cm}>{\raggedright\arraybackslash}p{4.7cm}@{}}
\toprule\noalign{}
日期 & 公众号 & 标题 & 可用信息 \\
\midrule\noalign{}
\endhead
\bottomrule\noalign{}
\endlastfoot
近一年 & 朱哥讲支付 & 一文读懂跨境支付业务详解 & 跨境支付、结汇和资金路径 \\
近一年 & 朱哥讲支付 & 外贸收款？关键是跨境结汇！ & 外贸收款与结汇流程 \\
近一年 & 朱哥讲金融 & 蚂蚁集团收购香港券商获批，正式切入持牌证券业务领域 &
香港持牌证券与数字金融边界 \\
\end{longtable}
}

\section{钱包、支付系统与运营}\label{ux94b1ux5305ux652fux4ed8ux7cfbux7edfux4e0eux8fd0ux8425}

{\def\LTcaptype{none} % do not increment counter
\begin{longtable}[]{@{}>{\raggedright\arraybackslash}p{1.9cm}>{\raggedright\arraybackslash}p{2.8cm}>{\raggedright\arraybackslash}p{6.0cm}>{\raggedright\arraybackslash}p{4.7cm}@{}}
\toprule\noalign{}
日期 & 公众号 & 标题 & 可用信息 \\
\midrule\noalign{}
\endhead
\bottomrule\noalign{}
\endlastfoot
近一年 & 朱哥讲支付 & 聚合支付的电子钱包设计 & 钱包账户、余额、场景和权限设计 \\
近一年 & 朱哥讲支付 & 聚合支付的代付业务详解 & 代付链路、资金与风控 \\
近一年 & 朱哥讲支付 & 聚合支付的代收业务详解 & 代收链路、商户与对账 \\
近一年 & 朱哥讲支付 & 聚合支付的对账体系设计 & 对账、差异和差错处理 \\
近一年 & 朱哥讲支付 & 支付路由系统设计 & 通道选择、成本、成功率与风控 \\
近一年 & 朱哥讲支付 & 支付对账系统设计 & 支付账务闭环 \\
近一年 & 朱哥讲支付 & 从支付角度回答``多级账户'' & 账户体系与资金层级 \\
近一年 & 朱哥讲支付 & 3500字，搞懂互联网支付 & 互联网支付基础知识 \\
近一年 & 朱哥讲支付 & 微信刷掌支付，如期而来 & 生物识别支付产品 \\
\end{longtable}
}

\section{作者书籍中的相关结构}\label{ux4f5cux8005ux4e66ux7c4dux4e2dux7684ux76f8ux5173ux7ed3ux6784}

《金融产品方法论》的公开目录/摘要提到``区块链金融 3
个运行模式''\,``区块链业务模式''和``区块链创新模式''。由于未获得可验证全文，本库只把它作为作者知识体系的线索，不扩写为具体结论。
